\scanpage{106}
\mspage{52}
11) Digression : unfolding of a divisor with normal crossings.

To fix ideas, I shall place myself in the context of smooth analytic multiplicities — but one could \unsure{also|see}, by taking \uline{relative} smooth analytic multiplicities — or by replacing the analytic assumption\unsure{s} by ad-hoc topological assumptions — or by passing to the scheme-theoretic context.....

Let $X$ be a smooth analytic multiplicity, and $D$ a divisor on $X$ with normal crossings (purely local assumption on $(X,D)$, by definition). \deleted{Let $G$ be a}

For every point $x$ of $D$, the set of ``branches'' of $D$ passing through $x$ is a finite set. $I(x)$, By definition, the \uline{germs} \inserted{at $x$} \uline{of sub-multiplicities} \uuline{attach\unsure{ed|ed}} \uuline{to $D$} are the intersections of a finite number \deleted{of} among the $D_{x,i}\in I(x)$.
\margin{it is better to take an $x\in X$, and set $I(x)=\phi$ if $x\notin D$.}
These \unsure{germs} form an ordered set \inserted{$\Delta_x\unsure{=(|(=}\Delta_x(X\unsure{,D}))$} anti-isomorphic \unsure{to|:} $\mathcal P(I(x))$, \deleted{If \unsure{$J$ is}}\inserted{the codimension} \deleted{\unsure{a|with|those}} of $D_{x,J}$ $(J\subset I(x))$ being given by
\[
(149)\qquad \operatorname{Codim}(D_{x,J},X_x)=\operatorname{card}J
\]
\margin{105}
